Table \ref{matrix_xg_test} shows the confusion matrices of XGBoost models, four with a single feature and one with all features. 
% 
Table \ref{matrix_xg_test_a} shows the confusion matrix of the model with image-based 50th percentile values as input features. 
% 
The model showed a trend of low FNR and low FPR.
% 
Table \ref{matrix_xg_test_b} shows the confusion matrix of the model with image-based 90th percentile values as input features. 
% 
The model showed a trend of low FNR and high FPR.
%
Table \ref{matrix_xg_test_c} shows the confusion matrix of the model with centroid 50th percentile values as input features. 
% 
The model showed a trend of medium FNR and low FPR. 
%
Table \ref{matrix_xg_test_d} shows the confusion matrix of the model with centroid 90th percentile values as input features. 
% 
The model showed a trend of high FNR and low FPR.
% 
Table \ref{matrix_xg_test_e} shows the confusion matrix of the model with all feature values as input. 
% 
The model showed a trend of low FNR and high FPR.
% 
Table \ref{result_xg_test} shows the values of the evaluation index for each XGBoost model.
% 
The model with the lowest FNR was the one that used both image-based 90th and 90th percentile values and all features.
%
The model with the lowest FPR was the one using centroid 50th percentile values.
% 
The model with the highest accuracy was the one using centroid 50th percentile values.
